% Standalone problem document, generated from the adjacent README.md.
% Compile with XeLaTeX. Mathematical content is maintained in Markdown.
\documentclass[11pt,a4paper]{article}
\usepackage[fontsize=10.5pt]{scrextend}
\usepackage[margin=22mm,headheight=15pt,headsep=9mm,footskip=11mm]{geometry}
\usepackage{fontspec}
\setmainfont{texgyrepagella}[Extension=.otf,UprightFont=*-regular,BoldFont=*-bold,ItalicFont=*-italic,BoldItalicFont=*-bolditalic]
\setsansfont{texgyreheros}[Extension=.otf,UprightFont=*-regular,BoldFont=*-bold,ItalicFont=*-italic,BoldItalicFont=*-bolditalic]
\setmonofont{texgyrecursor}[Extension=.otf,UprightFont=*-regular,BoldFont=*-bold,ItalicFont=*-italic,BoldItalicFont=*-bolditalic,Scale=MatchLowercase]
\usepackage{amsmath,amssymb}
\usepackage{unicode-math}
\setmathfont{texgyrepagella-math.otf}
\usepackage{xcolor,microtype,enumitem,needspace,fancyhdr}
\usepackage{bookmark}
\definecolor{ink}{HTML}{17344B}
\definecolor{accent}{HTML}{197D80}
\definecolor{muted}{HTML}{596773}
\hypersetup{colorlinks=true,linkcolor=accent,urlcolor=accent,pdftitle={AC-11 - The
least positive permanent of a sign
matrix},pdfauthor={Open Problems in Numerical Linear Algebra}}
\urlstyle{same}
\setlength{\parindent}{0pt}
\setlength{\parskip}{5pt plus 1pt minus 1pt}
\linespread{1.04}
\setlength{\emergencystretch}{3em}
\widowpenalty=10000
\clubpenalty=10000
\displaywidowpenalty=10000
\allowdisplaybreaks[1]
\setlist{nosep,leftmargin=1.5em}
\providecommand{\tightlist}{\setlength{\itemsep}{0pt}\setlength{\parskip}{0pt}}
\providecommand{\pandocbounded}[1]{#1}
\pagestyle{fancy}
\fancyhf{}
\fancyhead[L]{\sffamily\footnotesize\color{muted}OPEN PROBLEMS IN NUMERICAL LINEAR ALGEBRA}
\fancyhead[R]{\sffamily\bfseries\footnotesize\color{accent}AC-11}
\fancyfoot[L]{\parbox[b]{0.9\textwidth}{\sffamily\fontsize{8}{10}\selectfont\color{muted}Editorial ratings; a bounded search cannot certify that a problem remains unsolved.\\Literature check: 2026-09-11}}
\fancyfoot[R]{\sffamily\footnotesize\color{muted}\thepage}
\renewcommand{\headrulewidth}{0.3pt}
\renewcommand{\headrule}{\hbox to\headwidth{\color{accent}\leaders\hrule height \headrulewidth\hfill}}
\makeatletter
\renewcommand{\section}{\Needspace{5\baselineskip}\@startsection{section}{1}{0pt}{12pt plus 2pt minus 2pt}{4pt}{\sffamily\bfseries\large\color{ink}}}
\renewcommand{\subsection}{\Needspace{4\baselineskip}\@startsection{subsection}{2}{0pt}{9pt plus 1pt minus 1pt}{3pt}{\sffamily\bfseries\normalsize\color{ink}}}
\makeatother
\setcounter{secnumdepth}{0}
\begin{document}
{\sffamily\small\color{muted}Arithmetic and complexity\par}
\vspace{3pt}
{\raggedright\hyphenpenalty=10000\exhyphenpenalty=10000\sffamily\bfseries\fontsize{21}{24}\selectfont\color{ink}The
least positive permanent of a sign matrix\par}
\vspace{7pt}
{\sffamily\small\textbf{Difficulty:} challenging\quad\textbf{Importance:} interesting
to specialist\par}
{\sffamily\small\bfseries\color{accent}Status: Partially resolved\par}
\vspace{4pt}
{\color{accent}\hrule height 0.8pt}
\vspace{6pt}
\textbf{Rating rationale:} Challenging because divisibility alone gives
no construction attaining the bound for every order; specialist
importance concerns exact permanent values and discrete matrix
constructions.

\subsection{Independently reviewed finite cases -
2026-09-11}\label{independently-reviewed-finite-cases---2026-09-11}

Complete exact certificates attain \(2^{n-\lfloor\log_2(n+1)\rfloor}\)
for every order \(1\leq n\leq35\). Together with the proved universal
divisibility bound, they establish the conjectured minimum at these
orders, extending the cited order-20 construction range. The entire
supplied matrices are checked with two coprime moduli and a proved
uniqueness bound; cofactor dot products alone are not the evidence.
Attainment for every order \(n\geq36\) remains unresolved by this
submission.

\textbf{Author:} Matthew J. Colbrook, Department of Applied Mathematics
and Theoretical Physics, University of Cambridge.
\href{https://github.com/ajt60gaibb/OpenProblemsInNLA/blob/main/references/colbrook-arithmetic-2026-09-11/manuscripts/AC-11-12.pdf}{Complete
manuscript},
\href{https://github.com/ajt60gaibb/OpenProblemsInNLA/blob/main/references/colbrook-arithmetic-2026-09-11/verification/reviews/AC-11-12-review.md}{independent
review}, and
\href{https://github.com/ajt60gaibb/OpenProblemsInNLA/blob/main/references/colbrook-arithmetic-2026-09-11/README.md}{submission
record}. AI assistance is disclosed. Agent verification is not external
human peer review or formal certification; no novelty or priority claim
is made.

The
\href{https://github.com/ajt60gaibb/OpenProblemsInNLA/blob/main/references/colbrook-arithmetic-2026-09-11/verification/fresh/README.md}{fresh
certificate checks} complement the analytic review. The remaining
universal target and its ratings are retained.

\subsection{Problem statement}\label{problem-statement}

For \(n\ge1\), let \(\Omega_n=\{-1,1\}^{n\times n}\) and define \[
\operatorname{per}A=\sum_{\sigma\in S_n}\prod_{i=1}^n a_{i,\sigma(i)},
\qquad
p_n=\min\{\operatorname{per}A:A\in\Omega_n,\ \operatorname{per}A>0\}.
\] The set in this minimum is nonempty, since the all-ones matrix has
permanent \(n!\). Is Kräuter's proposed formula \[
p_n=2^{\,n-\lfloor\log_2(n+1)\rfloor}
\] valid for every \(n\)? The problem asks for attainability of the
divisibility lower bound, with no rank restriction on \(A\).

\subsection{Relevance and ratings}\label{relevance-and-ratings}

This is an extremal question about a basic multilinear matrix function
on discrete matrices. Exact values and divisibility constraints matter
for permanent computation and constructions of sign matrices. The
attainability problem across all orders merits challenging difficulty;
its main audience is combinatorial matrix theory.

\subsection{References}\label{references}

\begin{itemize}
\tightlist
\item
  I. M. Wanless,
  \href{https://users.monash.edu.au/~iwanless/papers/wangconjLAMA.pdf}{\emph{Permanents
  of matrices of signed ones}}, Linear and Multilinear Algebra
  \textbf{53} (2005), 427--433, §3, for attaining constructions through
  order twenty.
\item
  D. Ingram and A. Razborov, \emph{On the range of the permanent of
  \((\pm1)\)-matrices}, Linear Algebra Appl. 743 (2026), 271--285. The
  precise question is §6, Problem 3 in
  \href{https://arxiv.org/html/2507.09433v1}{arXiv:2507.09433v1}; the
  \href{https://doi.org/10.1016/j.laa.2026.04.027}{published
  introduction} gives the displayed unified formula.
\item
  M. V. Budrevich and A. E. Guterman, \emph{Kräuter conjecture on
  permanents is true}, J. Combin. Theory A 162 (2019), 306--343
  (\href{https://arxiv.org/abs/1810.04439}{paper};
  \href{https://doi.org/10.1016/j.jcta.2018.11.009}{journal}), for a
  different, resolved rank-dependent upper-bound conjecture.
\end{itemize}

\subsection{Status check --- 2026-09-10}\label{status-check-2026-09-10}

Checked
\href{https://doi.org/10.1016/j.laa.2026.04.027}{Ingram--Razborov's
August 2026 published introduction}, which explicitly retains the
minimum-positive-value conjecture, and its
\href{https://arxiv.org/html/2507.09433v1}{preprint §6, Problem 3}.
\href{https://users.monash.edu.au/~iwanless/papers/wangconjLAMA.pdf}{Wanless
(2005), §3, pp.~430--431}, supplies attaining matrices for every
\(n\le20\), a proved parameter range of this target. Exact-title and
minimum-positive-permanent searches found no general resolution. The
\href{https://arxiv.org/abs/1810.04439}{2019 Kräuter theorem} proves a
rank-dependent upper bound;
\href{https://arxiv.org/abs/2509.22577}{Hunter--Kwan--Sauermann} proves
exponential range cardinality. Neither establishes general attainment of
the lower bound.
\end{document}
